\documentclass[10pt, letterpaper]{article}
\usepackage[utf8]{inputenc}
\usepackage{geometry}
    \geometry{
        left=0.5in,
        right=0.5in,
        top=0.5in,
        bottom=0.75in,
    }
\usepackage{multicol}      
\usepackage{graphicx}
\usepackage{titlesec}        
\usepackage{booktabs}
\usepackage{palatino}
\usepackage{mathpazo}

\usepackage{caption}   
\usepackage{setspace}

\raggedbottom

\setstretch{1.0}  % Slight line spacing increase for readability
\setlength{\columnsep}{0.2in}  % Adjust column gap

\titleformat{\section}{\fontsize{10}{13}\bfseries}{\thesection.}{0.4em}{}
\titleformat{\subsection}{\fontsize{10}{13}\itshape}{\thesubsection.}{0.4em}{}

\titlespacing*{\section}{0pt}{1.5ex}{0.5em}
\titlespacing*{\subsection}{0pt}{1.25ex}{0.5em}

\newenvironment{mabstract}
    {
        \vspace{-0.5em}
        \noindent\rule{\linewidth}{1pt} % Top Line
        \par\vspace{0.3em}
        \noindent\textbf{Abstract} % Title (Left Aligned)
        \par\noindent\ignorespaces
    }
    {
        \par\vspace{0.2em}
        \noindent\rule{\linewidth}{1pt} % Bottom Line
        \vspace{-1.5em}
    }
    
\begin{document}

\begin{flushleft}
    {\LARGE\bfseries 
    Engine Fault Detection Using Machine Learning
    \par}
    \vspace{1.2em}
    {
    Qamer Abdo  - 1\textsuperscript{a}, Ali Ibrahim - 2\textsuperscript{b}, Faiyaz Chowdhury - 3\textsuperscript{a}, Araf Zaman - 4\textsuperscript{a}, Jakub Hann - 5\textsuperscript{C}
    \par}
    \vspace{1em}
    {\small\itshape
    \textsuperscript{a} Computer Science Major\\
    \textsuperscript{b} Mechanical Engineering Major\\
    \textsuperscript{c} Civil Engineering Major\\
    \par}
\end{flushleft}

\begin{mabstract}
An engine is central to an individual's ability to complete daily tasks, and an engine's failure would impose significant opportunity costs, maintenance costs, and safety incidents. Thus, this project utilizes the application of data-driven techniques and machine learning for predictive fault detection, with the goal of building a classifier capable of providing early warning of engine faults before they become critical. Unlike traditional threshold-based methods, machine learning is better suited for this problem as it can detect complex, multi-dimensional patterns across several sensor channels simultaneously. The dataset used consists of multiple sensor readings recorded over a series of time intervals, simulating sensor readings from an engine in an automotive application. It captures data related to engine performance, fault conditions, and operational modes over a series of time intervals. The engine readings are taken across different operating conditions. Comprising 55,999 observations, the dataset lists readings from fourteen sensors including RPM, speed, fuel efficiency, air fuel ratio, torque, power output, and emission measurements. Each labeled with one of four fault condition classes, with zero indicating normal operation and one through three representing distinct fault conditions. This report serves as an analysis of the process of this project covering data exploration, statistical characterization, and a review of the results and findings of our analysis.
\end{mabstract}

% --- BEGIN COLUMNS ---
\begin{multicols}{2}

\section{Introduction}
Engines are widely used around the world for things like transportation, manufacturing, and generating power, among other things. Because they are used in so many different ways, these systems are vulnerable to unexpected failures as they age. As a result, when these systems fail, it can lead to issues such as unplanned downtime, expensive repairs, and possible safety problems. For these reasons, it is important for engineers to come up with a solution to try and limit the risk of engine failure before it happens. 
Identifying potential failures in engines can be difficult due to the fact that most engine failure occurs slowly and may not be easy to recognize. In many cases, by the time engine failure is detected, significant damage may have already occurred. As technology improves, modern engines are equipped with different sensors to monitor engine parameters such as temperature, crankshaft rotation speed (rpm), pressure, and vibration. Through these sensor engineers are able to collect data about the engines condition and performance. 
Using this data, techniques such as machine learning can be used to analyze patterns based on the engines performance and predict failure before it happens. In this study, our main goal is to analyze engine sensor data, and code a machine learning model that will be able to recognize these patterns and detect if engine failure will occur or not based on different situations.

\begin{itemize}
    \item Objective 1: To analyze engine sensor data and identify patterns that would lead to potential engine failure.
    \item Objective 2: To implement a machine learning algorithm that is able to use the collected data and detect engine failure before it happens
\end{itemize}

\section{Literature Review}
For many years, engineers have used threshold-based methods to find engine problems.
This basically means setting fixed limits for sensor readings and sending an alert if a value goes
past that limit. While this method is easy to set up, it has an important weakness. It usually looks
at one sensor at a time, which means it can miss bigger problems that involve several sensors
working together. Rajapaksha et al. [5] explain that these traditional methods also depend a lot
on the experience of mechanics and special diagnostic tools, which can make the process
expensive and slow. As engines have become more advanced, this method is no longer enough
because modern engines create large amounts of data that these simple systems cannot handle
well [7].
This is where machine learning starts to help. In recent years, data-based methods have
become a common way to predict maintenance issues before they happen. Ersöz et al. [2] show
that models such as random forests, decision trees, and boosting methods can be trained using
sensor data. These models learn patterns from the data so they can spot unusual behavior, predict
failures, and help schedule maintenance earlier. Fernandes et al. [3] reviewed many of these
approaches across different industries and found that supervised classification is one of the most
widely used methods. In this type of model, the system learns from labeled examples so it can
later recognize when a fault is happening.
Researchers have also studied which machine learning models work best for this type of
problem. Acharya et al. [4] found that XGBoost often performs very well when predicting
equipment failures. They also found that deeper models, such as Long Short-Term Memory
networks, can sometimes reach even higher accuracy. When dealing with sensor data that
changes over time, which is the type of data used in this project, researchers have tested many
models. Some studies show that even simpler models like Logistic Regression can still perform
well when the data is prepared properly [8].
One major advantage of machine learning compared to older methods is its ability to find
patterns across many sensors at once. Chen et al. [9] explain that a single sensor reading may not
tell the full story. The meaning of that reading can change depending on what other sensors are
showing at the same time. A simple threshold rule would miss these relationships completely.
Bucur et al. [6] expand on this idea for engine diagnostics. Their work shows that combining
different measurements such as compression, fuel angle, and vibration into one dataset allows
models to better judge the overall condition of an engine over time.
When looking specifically at internal combustion engines, many studies support the use
of machine learning for fault detection. Rajapaksha et al. [5] note that neural networks work well
in this area because they can learn the connections between sensor readings and engine faults
without engineers having to manually define every rule. Other models such as decision trees,
random forests, and support vector machines have also shown strong results. Kumar et al. [10]
demonstrate how these models can be used to predict equipment failures in real systems.
One challenge that researchers often face is that normal engine readings appear much
more often than fault readings in real datasets. Because of this, models may learn to focus too
much on normal cases and ignore the rare fault cases. Zhang et al. [11] found that models like
support vector machines and standard neural networks sometimes struggle when the data is
unbalanced in this way. Newer methods, such as mixed sampling techniques and bidirectional
LSTM networks, have shown better performance when dealing with this problem [11]. Recently, Vergara et al. conducted a study in which they used sensors on a C14NE spark ignition engine to collect data on how the engine performs under different simulated conditions [12]. The data that they collected was made specifically for the implementation of machine learning software which is we selected it for use in this project. 
Overall, past research strongly supports using machine learning for engine fault
detection. The goal of this project is to apply these ideas to a structured dataset that includes
 different simulated engine conditions. By training models on this data,
the project aims to build a system that can identify possible faults early before they turn into
more serious engine problems.

\section{Data}

    The dataset used in this study comes from the EngineFaultDB dataset introduced by Vergara et al. in the study, the data was collected through the use of various different sensors to monitor 14 different engine parameters as seen in table 1. The dataset contains 55,999 observation where each observation represents the engines condition at a given time. The data was checked in MATLAB using the "summary" command. This tool was useful because not only did it calculate the mean and standard deviation of each sensor, it also checked to see if the data was missing any values. Since there were no missing values detected, there was no further preprocessing that needed to be done. 
    The dataset includes several sensor measurements to analyze the engines operating condition. These variables include Manifold absolute pressure (MAP), Throttle position sensor (TPS), Force, Power, Crankshaft rotation (RPM), Fuel consumption, Speed, Air-fuel ratio, Lambda, Oxygen (O$_2$), Hydrocarbon (HC), Carbon monoxide (CO), and Carbon Dioxide (CO$_2$).

% --- SINGLE COLUMN TABLE ---
% Note: We use 'center' environment + \captionof instead of 'table' environment
\begin{center}
    \captionof{table}{Data Summary Statistics}
    \label{tab:data_summary}
    \begin{tabular}{lcc}
        \toprule
        \textbf{Variable} & \textbf{Mean} & \textbf{Std. Dev}\\
        \midrule
MAP & 1.8325 & 0.8378 \\
TPS & 1.3953 & 0.9070 \\
Force & 286.6917 & 378.7749 \\
Power & 5.6571 & 7.6841 \\
RPM & 2398.1 & 932.0087 \\
Consumption$_{L/H}$ & 4.4976 & 2.2215 \\
Consumption$_{L/100km}$ & 8.9398 & 3.1541 \\
Speed & 51.6887 & 20.1403 \\
AFR & 14.1694 & 0.9698 \\
Lambda & 0.9639 & 0.0660 \\
O$_2$ & 0.5859 & 0.2237 \\
HC & 188.4461 & 111.0470 \\
CO & 1.9324 & 1.9888 \\
CO$_2$ & 13.0352 & 1.0466 \\
        \bottomrule
    \end{tabular}
\end{center}

\end{multicols}

% --- BREAK COLUMNS FOR WIDE IMAGE ---
\vfill
\noindent\makebox[\textwidth][c]{%
    \includegraphics[width=0.30\textwidth,height=1.3in,keepaspectratio]{fault_histogram}\hfill%
    \includegraphics[width=0.30\textwidth,height=1.3in,keepaspectratio]{fault_barchart}\hfill%
    \includegraphics[width=0.30\textwidth,height=1.3in,keepaspectratio]{fault_piechart}%
}
\par\vspace{0.2em}
\captionof{figure}{Fault distribution analysis: histogram (left), bar chart (centre), and pie chart (right).}
\label{fig:wide_results}
\vfill

\newpage
\begin{multicols}{2}

\section{Methodology}
This section details \textit{how} you conducted your research. It should be rigorous enough that another student could replicate your work.

If you used a specific formula, ensure it is formatted correctly:
\begin{equation}
    E = mc^2
    \label{eq:einstein}
\end{equation}

Describe your experimental setup, algorithms used, or the qualitative analysis process.

\section{Results}
Present your findings here objectively. As seen in Figure \ref{fig:wide_results}, there is a significant correlation between the variables. Discuss what the numbers mean in the context of your initial problem.

You can also include smaller figures inside the column if needed:

\begin{center}
    \framebox{\parbox{\linewidth}{\centering \vspace{1cm} [Small Figure] \vspace{1cm}}}
    \captionof{figure}{A column-width detailed view.}
    \label{fig:small_detail}
\end{center}

\section{Conclusion}
Summarize your main findings and their implications. Discuss the limitations of your study (e.g., small sample size, limited time) and suggest directions for future work.


\begin{thebibliography}{9}
\bibitem{theissler2021}
A. Theissler, J. Pérez-Velázquez, M. Kettelgerdes, and G. Elger,
``Predictive maintenance enabled by machine learning: Use cases and challenges in the automotive industry,''
\textit{Reliability Engineering \& System Safety}, vol. 215, 107864, 2021.

\bibitem{ersoz2022}
O. Ersöz, A. İnal, A. Aktepe, A. Türker, and S. Ersöz,
``A systematic literature review of predictive maintenance from a transportation systems aspect,''
\textit{Sustainability}, vol. 14, no. 21, 14536, 2022.

\bibitem{fernandes2022}
M. Fernandes, J. M. Corchado, and G. Marreiros,
``Machine learning techniques applied to mechanical fault diagnosis and fault prognosis in the context of real industrial manufacturing use-cases,''
\textit{Applied Intelligence}, vol. 52, pp. 14246--14280, 2022.

\bibitem{acharya2024}
V. Acharya et al.,
``Predicting machine failures using machine learning and deep learning algorithms,''
\textit{Results in Control and Optimization}, 2024.

\bibitem{rajapaksha2024}
N. Rajapaksha et al.,
``Machine learning approaches for fault detection in internal combustion engines,''
\textit{Informatics}, vol. 12, no. 1, 25, 2024.

\bibitem{bucur2023}
A. Bucur et al.,
``Using event data to build predictive engine failure models,''
\textit{Machines}, vol. 11, no. 7, 704, 2023.

\bibitem{abdul2025}
Z. Abdul et al.,
``Research progress on data-driven industrial fault diagnosis methods,''
\textit{PMC}, 2025.

\bibitem{leukel2024}
J. Leukel et al.,
``Predicting machine failures from multivariate time series: an industrial case study,''
\textit{arXiv}, 2024.

\bibitem{chen2025}
H. Chen et al.,
``Research on bearing fault diagnosis based on machine learning and SHAP interpretability analysis,''
\textit{Scientific Reports}, 2025.

\bibitem{kumar2023}
A. Kumar et al.,
``Predictive maintenance of armoured vehicles using machine learning approaches,''
\textit{arXiv}, 2023.

\bibitem{zhang2025}
W. Zhang et al.,
``Fault prediction of aircraft engine based on adaptive hybrid sampling and BiLSTM,''
\textit{Scientific Reports}, 2025.

\bibitem{vergara2023}
M. Vergara, L. Ramos, N. D. Rivera-Campoverde, and F. Rivas-Echeverría,
``EngineFaultDB: A Novel Dataset for Automotive Engine Fault Classification and Baseline Results,''
\textit{IEEE Access}, vol. 11, pp. 126155--126171, 2023,
doi: 10.1109/ACCESS.2023.3331316.

\end{thebibliography}

\end{multicols}

\end{document}